% eq_calibration_constants.tex — Three Initial Hardware Boundaries (Calibration Constants)
% Single source of truth. \input this wherever the macroscopic calibration scale
% is formally restated.
%
% Per Vol 1 Ch 1's canonical structure (manuscript/vol_1_foundations/chapters/01_fundamental_axioms.tex
% lines 14-21 + 218-222): the macroscopic continuum dynamics rests on FOUR
% foundational axioms (eq_axiom_{1,2,3,4}.tex). The numerical calibration
% boundaries below are DERIVED from those axioms; they set the empirical scale
% at which the axioms operate.
%
% Axiom FORM-attributions (values are calibration inputs, per FORM/VALUE signature):
%   Z_0      <- Axiom 1 (LC network, sqrt(mu_0/epsilon_0))
%   ell_node <- Axiom 1 (unknot ground-state ropelength)
%   alpha    <- Class-B named identification; value = STANDING echo (3 routes closed)
%   xi_topo  <- Axiom 2 ([Q] equiv [L])
%   V_snap   <- Axiom 4, longitudinal A1 dilatation (rest energy m_e c^2; def-vyvsn1)
%   V_yield  <- Axiom 4, transverse Cosserat T2 self-trap wall (sqrt(alpha)*V_snap;
%               def-vyvsn1 = T2, PR#260; sqrt(alpha) is an alpha-echo, not independent)
%
% Newton's constant G is a fourth derived calibration parameter, emerging from
% Machian thermodynamic equilibrium of the LC lattice. Its formal definition and
% the gravitational refractive-index derivation it generates appear in
% eq_gravity_derived.tex (separated because gravity is a derived consequence of
% Ax 1 + Ax 4 under Symmetric Scaling, not a primitive axiom).
%
\begin{resultbox}{Master Calibration Constants}
The numerical scale of the macroscopic continuum is set by the following constants.
Their \textbf{FORM} is attributed to the four foundational axioms (cited inline); their
\textbf{VALUES} are \textbf{calibration inputs the substrate does not independently
select} --- the framework's FORM-deriving / VALUE-importing signature (KB leaf
\texttt{form-deriving-value-importing.md}). ``Derived'' below means form-attributed,
not value-forced.

\medskip

\noindent\textbf{Lattice Pitch} (per Axiom 1, unknot ropelength):
\begin{equation}
    \ell_{node} = \frac{\hbar}{m_e c} \approx 3.86 \times 10^{-13}\;\text{m}
    \label{eq:cal_lnode}
\end{equation}

\noindent\textbf{Vacuum Impedance} (per Axiom 1, LC constituents):
\begin{equation}
    Z_0 = \sqrt{\frac{\mu_0}{\varepsilon_0}} \approx 376.73\;\Omega
    \label{eq:cal_z0}
\end{equation}
where $\mu_0$ is the per-node inductance (rotational inertia) and
$\varepsilon_0$ is the per-node capacitance (elastic compliance).

\noindent\textbf{Topological Conversion Constant} (per Axiom 2, $[Q]\equiv[L]$):
\begin{equation}
    \xi_{topo} \equiv \frac{e}{\ell_{node}} \approx 4.149 \times 10^{-7}\;\text{C/m}
    \label{eq:cal_xi_topo}
\end{equation}

\noindent\textbf{Fine-Structure Constant} (Class-B named identification; value = \textbf{STANDING echo}):
\begin{equation}
    \alpha = \frac{e^2}{4\pi\varepsilon_0\hbar c} = \frac{p_c}{8\pi} \approx \frac{1}{137.036}, \qquad
    \alpha^{-1} = 4\pi^3 + \pi^2 + \pi
    \label{eq:cal_alpha}
\end{equation}
The closed-form $\alpha^{-1} = 4\pi^3 + \pi^2 + \pi$ is \textbf{Class-B substrate-mechanism manifestation via a named identification} ($R\cdot r = 1/4$ on the Golden Torus) that the substrate does \emph{not} independently select. The value is a \textbf{STANDING echo} (Grant-ratified 2026-06-18): all three named lift-routes are \textbf{closed-negative} --- dynamical selection (flat), kinematic unit-bridge (Class-B, absorbs $\alpha$), and the rigidity-percolation count (structurally closed at the central-force isostatic ceiling $z_{eff} \to 6$, \texttt{alpha\_free\_map\_to\_137\_exists = False}). NOT a first-principles derivation and NOT a zero-parameter closure; the flip-condition ($R\cdot r = 1/4$ forced without $\alpha$-circularity) stays live. Canon: KB \texttt{ch8-alpha-golden-torus.md}, \texttt{form-deriving-value-importing.md}.

\noindent\textbf{Dielectric Snap Voltage} (per Axiom 4; the \textbf{longitudinal A1 dilatation} saturation / rest-energy scale, $m_e c^2$; \texttt{def-vyvsn1}, PR\,\#260):
\begin{equation}
    V_{snap} = \frac{m_e c^2}{e} \approx 511.0\;\text{kV}
    \label{eq:cal_vsnap}
\end{equation}

\noindent\textbf{Dielectric Yield Voltage} (per Axiom 4; the \textbf{transverse Cosserat $T_2$ self-trap wall} --- the electron's confining $\Gamma = -1$; \texttt{def-vyvsn1} adjudicated = $T_2$, PR\,\#260):
\begin{equation}
    V_{yield} = \sqrt{\alpha} \cdot V_{snap} \approx 43.65\;\text{kV}
    \label{eq:cal_vyield}
\end{equation}
Because $V_{yield} = \sqrt{\alpha}\, V_{snap}$ \emph{exactly}, the two are \textbf{not} independent per-sector thresholds: the $\sqrt{\alpha}$ is an \textbf{$\alpha$-echo} (Class-C), and the A1 mass core operates sub-saturated at $A = V_{yield}/V_{snap} = \sqrt{\alpha} \approx 0.085$ ($S \approx 0.996$) inside the $T_2$ cage --- which is why the mass channel binds without the $S \to 0$ runaway firing.
\end{resultbox}
